% !PRINTED 12
\weberpage{12}
%<<CONT
oder gleich oder grösser als $\mu$. Ist $e'$ gleich $\mu$, so liegen
sowohl zwischen $e$ und $e'$ als zwischen $e'$ und $\epsilon$ rationale
Verhältnisse. Ist aber $e'$ kleiner als $\mu$, so liegt $\mu$ zwischen
$e'$ und $\epsilon$, und ist $e'$ grösser als $\mu$, so liegt $\mu$
zwischen $e$ und $e'$; also kann $e'$ nicht zugleich gleich $e$ und
gleich $\epsilon$ sein.

Im Falle 2) heissen die Verhältnisse $e,\epsilon$ \gesperrt{ungleich}.
Es kann also entweder $\alpha)$ $e<\mu<\epsilon$ oder $\beta)$
$e>\mu'>\epsilon$ sein, und diese beiden Fälle schliessen sich aus,
weil aus $e<\mu<\epsilon$, $\epsilon<\mu'$ folgt, dass $\mu<\mu'$,
und folglich $e<\mu'$ sich ergiebt.

Die Grössenbeziehung $2\alpha)$ oder $2\beta)$ bleibt auch bestehen,
wenn $e$ oder $\epsilon$ durch ein ihm gleiches Element ersetzt wird.
Denn ist $\mu<\epsilon$ und $\mu\gtrless\epsilon'$, so liegt zwischen
$\epsilon$ und $\epsilon'$ ein rationales Verhältniss und
$\epsilon,\epsilon'$ sind nicht gleich.

Im Falle $2\alpha)$ heisst $e$ \gesperrt{kleiner als} $\epsilon$, im
Falle $2\beta)$ heisst $e$ \gesperrt{grösser als} $\epsilon$.

Zwischen zwei ungleichen Verhältnissen kann man eine beliebige Anzahl
rationaler Verhältnisse einschieben.

Wenn wir nun alle unter einander gleichen Verhältnisse zusammenfassen,
so erhalten wir einen Gattungsbegriff, den wir als \gesperrt{Zahl} im
allgemeinen Sinne des Wortes bezeichnen. Die Zahl ist also ein Name
oder Zeichen für eine gewisse Mannigfaltigkeit, deren Elemente eben
die mit einem unter ihnen gleichen Verhältnisse sind.\setcounter{footnote}{0}\footnote{Auf den
Gattungsbegriff lassen sich auch die natürlichen Zahlen in einfacher
und folgerichtiger Weise zurückführen.} Unter diesem Zahlbegriff sind
die rationalen Verhältnisse und folglich auch die natürlichen Zahlen
als die Verhältnisse $m:1$ mit enthalten und bilden die
\gesperrt{rationalen Zahlen}.

Zahlen, die nicht aus rationalen Verhältnissen entspringen, heissen
\gesperrt{irrationale Zahlen}.

Nach dem, was bis jetzt ausgeführt ist, bilden die Zahlen eine
\gesperrt{geordnete Menge}, und man kann ihre Ordnung feststellen,
wenn für jede Zahl irgend eines der darunter enthaltenen Verhältnisse
als Repräsentant gewählt wird.

\begin{center}
\rule{0.18\linewidth}{0.4pt}
\end{center}

Von zwei Verhältnissen mit demselben Nenner und ungleichen Zählern ist
das das grössere, dessen Zähler grösser ist, und
%CONT>>
